\begin{tomtat}
\begin{itemize}
  \item Với \tn{mô hình ngôn ngữ lớn}{large language model}, cần phân biệt đầu
  vào token, trạng thái do kiến trúc tạo ra,
  và hành vi đầu ra mà người dùng quan sát được.
  \item Transformer dùng \tn{cơ chế chú ý}{attention mechanism} để tạo biểu diễn theo ngữ cảnh, nhưng
  attention quan sát được không tự là bằng chứng nhân quả cho một token đầu ra.
  \item \tn{Tiền huấn luyện}{pretraining}, \tn{lời nhắc}{prompt} và RLHF đều có thể
  làm thay đổi hành vi cần được giải thích.
  \item \tn{Chuỗi suy luận}{chain-of-thought} là văn bản mô hình sinh ra dưới một
  lời nhắc đã chọn; sự tương ứng của nó với cơ chế nội tại vẫn phải được chứng
  minh riêng.
  \item Một cách đánh giá lời giải thích phải nêu rõ checkpoint, đầu vào và
  can thiệp, rồi mới diễn giải thay đổi ở đầu ra.
\end{itemize}
\end{tomtat}

\cauhoi
\begin{cauhoids}
  \item Một lời nhắc có hai ví dụ few-shot và một chỉ dẫn. Nếu thay một ví dụ
  làm câu trả lời đổi, phát biểu nào về hành vi đầu ra được hỗ trợ, và phát
  biểu nào về trọng số mô hình chưa được hỗ trợ?
  \item Đọc bài \enquote{Attention Is All You Need}~\autocite{p01attention}.
  Hãy mô tả \tn{tự chú ý}{self-attention} như một phép tính trên biểu diễn, rồi nêu can thiệp
  nào cần có trước khi gọi một trọng số attention là nguyên nhân của đầu ra.
  \item Đọc bài InstructGPT~\autocite{p04instructgpt}. Hãy chỉ ra dữ liệu nào
  đi vào tinh chỉnh có giám sát, dữ liệu nào đi vào \tn{mô hình thưởng}{reward model},
  và vì sao hai
  nguồn dữ liệu này không tạo thành lời giải thích cho từng câu trả lời.
  \item Một hệ thống trả lời đúng hơn khi được yêu cầu viết chuỗi suy luận.
  Dựa trên bài về chuỗi suy luận~\autocite{p06cot}, hãy phân biệt kết quả
  thực nghiệm đó với một kết luận về độ trung thực của chuỗi đã in ra.
\end{cauhoids}
